% =====================================================================
% 10-backup.tex -- Backup slides (not counted in the 60 minutes) and
% references. Owner: T3; references: E2.
% =====================================================================
\renewcommand{\sectioninfo}{Backup slides and references}
\section{Backup}

% ------------------------------------------------------------------
\begin{frame}{Backup: all pilot runs (seed 1, 100 steps)\hfill\ours}
  \centering\footnotesize
  \renewcommand{\arraystretch}{1.2}
  \begin{tabular}{@{}l l r r r@{}}
    \toprule
    \textbf{Run} & \textbf{Arm} & \textbf{lr} & \textbf{sampling $T$} & \textbf{Accuracy (\%)} \\
    \midrule
    \arm{pilot\_human\_s1}         & human         & $1\cdot10^{-5}$ & --  & 51.0 \\
    \arm{pilot\_teacher\_s1}       & teacher       & $1\cdot10^{-5}$ & --  & 54.4 \\
    \arm{pilot\_ref\_hardlabel\_s1} & ref\_hardlabel & $1\cdot10^{-5}$ & --  & 29.6 \\
    \midrule
    \arm{pilot\_onpolicy\_lr1e-05\_t0.7\_s1} & onpolicy & $1\cdot10^{-5}$ & 0.7 & 60.6 \\
    \arm{pilot\_onpolicy\_lr2e-05\_t0.7\_s1} & onpolicy & $2\cdot10^{-5}$ & 0.7 & 50.2 \\
    \arm{pilot\_onpolicy\_lr5e-06\_t0.5\_s1} & onpolicy & $5\cdot10^{-6}$ & 0.5 & 64.5 \\
    \arm{pilot\_onpolicy\_lr5e-06\_t0.7\_s1} & onpolicy & $5\cdot10^{-6}$ & 0.7 & \textbf{65.6} \\
    \arm{pilot\_onpolicy\_lr5e-06\_t1\_s1}   & onpolicy & $5\cdot10^{-6}$ & 1.0 & 63.5 \\
    \midrule
    \arm{ref\_student0} & untrained & -- & -- & 62.0 \\
    \bottomrule
  \end{tabular}

  \vspace{0.2cm}
  \raggedright\scriptsize
  At $1\cdot10^{-5}$ every off-policy pilot fell below the untrained student, which led to the
  sweep and to $5\cdot10^{-6}$ for all main runs. Pilots and main runs used different learning
  rates, so their numbers are not compared with each other. All 8 pilots saw identical
  questions in identical order.
\end{frame}

% ------------------------------------------------------------------
\begin{frame}[fragile]{Backup: running one arm on Kaggle\hfill\who{T3}}
  \begin{columns}[T]
    \begin{column}{0.58\textwidth}
\begin{lstlisting}[basicstyle=\ttfamily\tiny]
# notebooks/kaggle_run.ipynb, settings cell
REPO = "https://github.com/redkitex/Knowledge-Distillation-and-Model-Compression.git"
REF = "f2accc6"                # same commit for every seed
CONFIG = "configs/default.yaml"
MODE = "train+eval"            # or "eval" for references
ARM = "onpolicy"               # human | teacher | onpolicy | ref_hardlabel
SEED = 3
RUN_ID = None                  # -> "onpolicy_s3"
ACCOUNT = "T3"                 # whose Kaggle account
MAX_STEPS = None               # None = the full 500 steps
\end{lstlisting}
    \end{column}
    \begin{column}{0.39\textwidth}
      \footnotesize
      \begin{enumerate}
        \item Session options: GPU T4 x2, internet on.
        \item Edit the settings cell only.
        \item Save Version, Save \& Run All: clone, install, train, evaluate.
        \item Download \texttt{<run\_id>.zip} from the output; copy the row
              from \texttt{run\_info.json} into \repofile{results/run\_log.csv}.
      \end{enumerate}
      \vspace{0.1cm}
      {\scriptsize\color{kdgray} One free account runs one GPU session at a time:
      start the 3-hour on-policy run first.}
    \end{column}
  \end{columns}
\end{frame}

% ------------------------------------------------------------------
\begin{frame}{Backup: how the project ran\hfill\who{T3}}
  \centering
  \begin{tikzpicture}[font=\tiny,
      ev/.style={text width=1.35cm,align=center,inner sep=1pt},
      done/.style={circle,fill=kdblue,inner sep=1.4pt},
      plan/.style={circle,draw=kdgray,fill=white,inner sep=1.4pt}]
    \def\ux{0.72}
    \draw[kdgray,thick] (-0.3,0) -- (12.6,0);
    % day ticks every two days, labelled
    \foreach \dd/\lab in {0/15,2/17,4/19,6/21,8/23,10/25,12/27,14/29}
      {\draw[kdgray] (\dd*\ux,-0.08) -- (\dd*\ux,0.08);
       \node[text=kdgray,below=0.1cm] at (\dd*\ux,0) {Nov \lab};}
    \draw[kdgray] (16*\ux,-0.08) -- (16*\ux,0.08);
    \node[text=kdgray,below=0.1cm] at (16*\ux,0) {Dec 1};
    % pilots range
    \fill[kdblue!25] (6*\ux,-0.05) rectangle (9*\ux,0.05);
    % events above
    \node[done] at (0,0) {};             \node[ev,above=0.25cm] at (0,0) {plan restarted};
    \node[done] at (3*\ux,0) {};         \node[ev,above=0.25cm] at (3*\ux,0) {training, evaluation, stats code merged; Project Meeting 3};
    \node[ev,above=0.25cm] at (7.5*\ux,0) {pilots and learning-rate sweep (Nov 21 to 24)};
    \node[plan] at (14*\ux,0) {};        \node[ev,above=0.25cm,text=kdgray] at (14*\ux,0) {drafts done};
    \node[plan] at (16*\ux,0) {};        \node[ev,above=0.25cm,text=kdgray] at (16*\ux,0) {submit artifacts};
    % events below
    \node[done] at (1*\ux,0) {};         \node[ev,below=0.55cm] at (1*\ux,0) {skeleton and \texttt{data.py} merged};
    \node[done] at (4*\ux,0) {};         \node[ev,below=0.55cm] at (4*\ux,0) {2,000 teacher answers cached};
    \node[done] at (9.5*\ux,0) {};       \node[ev,below=0.55cm] at (9.5*\ux,0) {lr locked; seeds 2 and 3 trained and logged};
    \node[plan] at (15*\ux,0) {};        \node[ev,below=0.55cm,text=kdgray] at (15*\ux,0) {dress rehearsal};
    \node[plan,draw=kdorange] at (17*\ux,0) {}; \node[ev,below=0.55cm,text=kdorange] at (17*\ux,0) {\textbf{talk}, Dec 2};
    % today marker
    \draw[kdorange,dashed] (10*\ux,-0.45) -- (10*\ux,0.45);
  \end{tikzpicture}

  \vspace{0.35cm}
  \scriptsize
  \begin{tabular}{@{}llll@{}}
    \toprule
    & \textbf{1} & \textbf{2} & \textbf{3} \\
    \midrule
    \textbf{Training team} (builds, runs) & T1 training loop & T2 data, teacher answers & T3 runs, configs, slides \\
    \textbf{Evaluation team} (measures, explains) & E1 statistics, note & E2 evaluation, literature & E3 figures, repository \\
    \bottomrule
  \end{tabular}
  \vspace{0.1cm}

  {\scriptsize\color{kdgray} Filled dots: done (from the commit history); open dots: planned; dashed line: Nov 25.}
\end{frame}

% ------------------------------------------------------------------
\begin{frame}{References\hfill\who{E2}}
  \begin{columns}[T]
    \begin{column}{0.49\textwidth}
      \tiny
      \begin{itemize}\setlength{\itemsep}{1pt}
        \item Agarwal et al. (2024). On-Policy Distillation of Language Models: Learning from Self-Generated Mistakes. ICLR.
        \item Bengio, Vinyals, Jaitly, Shazeer (2015). Scheduled Sampling for Sequence Prediction with Recurrent Neural Networks. NeurIPS.
        \item Brown, Cai, DasGupta (2001). Interval Estimation for a Binomial Proportion. Statistical Science.
        \item Cobbe et al. (2021). Training Verifiers to Solve Math Word Problems. arXiv.
        \item DeepSeek-AI (2025). DeepSeek-R1: Incentivizing Reasoning Capability in LLMs via Reinforcement Learning. arXiv.
        \item Efron, Tibshirani (1993). An Introduction to the Bootstrap. Chapman \& Hall.
        \item Furlanello et al. (2018). Born Again Neural Networks. ICML.
        \item Gu, Dong, Wei, Huang (2024). MiniLLM: Knowledge Distillation of Large Language Models. ICLR.
        \item Hinton, Vinyals, Dean (2015). Distilling the Knowledge in a Neural Network. arXiv.
        \item Kim, Rush (2016). Sequence-Level Knowledge Distillation. EMNLP.
      \end{itemize}
    \end{column}
    \begin{column}{0.49\textwidth}
      \tiny
      \begin{itemize}\setlength{\itemsep}{1pt}
        \item Li et al. (2026). Rethinking On-Policy Distillation of Large Language Models: Phenomenology, Mechanism, and Recipe. arXiv.
        \item Lu, Thinking Machines Lab (2025). On-Policy Distillation. Blog post.
        \item Miller (2024). Adding Error Bars to Evals: A Statistical Approach to Language Model Evaluations. arXiv.
        \item Qwen Team (2025). Qwen3 Technical Report. arXiv.
        \item Ross, Gordon, Bagnell (2011). A Reduction of Imitation Learning and Structured Prediction to No-Regret Online Learning. AISTATS.
        \item Song, Zheng (2026). A Survey of On-Policy Distillation for Large Language Models. arXiv.
        \item Stanton et al. (2021). Does Knowledge Distillation Really Work? NeurIPS.
        \item Wilson (1927). Probable Inference, the Law of Succession, and Statistical Inference. JASA.
        \item Wu et al. (2025). Reasoning or Memorization? Unreliable Results of Reinforcement Learning Due to Data Contamination. arXiv.
      \end{itemize}
    \end{column}
  \end{columns}
  \vspace{0.1cm}
  \begin{todobox}{Verify every reference (E2)}
    The course requires checking authors, year and venue against each paper. Ross 2011, Li 2026 and
    Song and Zheng 2026 are not yet in \repofile{report/references.bib}.
  \end{todobox}
\end{frame}
